\input{sol_pre.tex}
\input{q_2025-02-23_2.tex}
\begin{document}
\input{sol_head.tex}

\Savol{1}
$11$ ga köpaytiramiz: $55a+5b=495$ $\Rightarrow$ $11a+b=99$. Raqamlar: $a=9$, $b=0$. \ $2a+b=18$.
\Javob{B}

\Savol{2}
$y$ ning har bir hadi $x$ ning mos hadidan $2$ marta katta: $y=2x$, \ $\dfrac yx=2$.
\Javob{A}

\Savol{3}
Unumdorliklar $p_1,p_2$ (töldiradi), $p_3,p_4$ (böşatadi):
\[p_1+p_2-p_3-p_4=\tfrac15,\quad p_1+p_2-p_3=\tfrac13\ \Rightarrow\ p_4=\tfrac2{15};\qquad p_1-p_3-p_4=\tfrac1{30}\ \Rightarrow\ p_1-p_3=\tfrac16.\]
Javob: $6$ soat.
\Javob{C}

\Savol{4}
Quruq modda özgarmaydi: $0{,}2x=0{,}4(x-25)$ $\Rightarrow$ $x=50$.
\Javob{C}

\Savol{5}
Surat: $3^3\cdot5^4+5^4\cdot3^3\cdot5=27\cdot625\cdot6$. \ Maxraj: $450-36\cdot\frac{25}4=225$.
\[\frac{27\cdot625\cdot6}{225}=450.\]
\Javob{C}

\Savol{6}
$a^2-14a+51=(a-7)^2+2=3+2=5$.
\Javob{B}

\Savol{7}
$\sqrt[4]{45}=\sqrt3\cdot\sqrt[4]5$: \ surat $3\left(1-\sqrt[4]5\right)^2$, maxraj $\left(1-\sqrt[4]5\right)^2$. \ Natija $3$.
\Javob{D}

\Savol{8}
$a_4+a_{14}=a_7+a_{11}=2a_9$ $\Rightarrow$ $5a_9=45$, $a_9=9$. \ $S_{17}=17a_9=153$.
\Javob{A}

\Savol{9}
$\dfrac{S_8}{S_4}=1+q^4=82$ $\Rightarrow$ $q^4=81$.\\
$q=3$: $S_4=40b_1=80$, $b_1=2$, \ $S_6=\dfrac{2\left(3^6-1\right)}2=728$. \ ($q=-3$ da ham $b_1=-4$, $S_6=728$.)
\Javob{C}

\Savol{10}
$a<c<0<b$: \ $|b-c|=b-c$, \ $|c-a|=c-a$, \ $|a-b|=b-a$; \ $\frac a{|a|}=\frac c{|c|}=-1$, $\frac b{|b|}=1$:
\[-(b-c)+(c-a)-(b-a)=2c-2b.\]
\Javob{D}

\Savol{11}
$\dfrac{a^4-64}{a^2+8}=a^2-8$, \ $\dfrac{\left(a^2\right)^3-8^3}{a^4+8a^2+64}=a^2-8$. \ Ayirma $0$.
\Javob{C}

\Savol{12}
$1-2\sin^2x-1=5\sin x+2$ $\Rightarrow$ $2\sin^2x+5\sin x+2=0$ $\Rightarrow$ $\sin x=-\frac12$.\\
Oraliq uzunligi $8\pi$ — $4$ ta töla davr, har birida $2$ ta ildiz (uçlarida $\sin x=1$): $8$ ta.
\Javob{C}

\Savol{13}
$\tg\left(\frac\pi4-\frac\alpha2\right)=\dfrac{1-\sin\alpha}{\cos\alpha}$:
\[\frac{\sin\alpha\cos\alpha}{(1-\sin\alpha)(1+\sin\alpha)}=\frac{\sin\alpha\cos\alpha}{\cos^2\alpha}=\tg\alpha.\]
\Javob{A}

\Savol{14}
$\left(x^2-1\right)\left(3^x-9\right)>0$, nollar $-1;\ 1;\ 2$: \ $x\in(-1;\,1)\cup(2;\,\infty)$.
\Javob{C}

\Savol{15}
$\log_3\log_2g\ge1$ $\Rightarrow$ $\log_2g\ge3$ $\Rightarrow$ $g\ge8$:
\[\frac{x^3+8x}{x-5}-8=\frac{x^3+40}{x-5}\ge0\ \Rightarrow\ x\le-\sqrt[3]{40}\ \text{yoki}\ x>5.\]
Eng kiçik natural yeçim $6$.
\Javob{D}

\Savol{16}
$x\left(245-x^5\right)>0$, $x\in\mathbb N$: $x^5<245$ $\Rightarrow$ $x=1,2,3$ ($3^5=243$). \ Yiğindi $6$.
\Javob{C}

\Savol{17}
Ayiramiz: $x^2-y^2=-(x-y)$ $\Rightarrow$ $(x-y)(x+y+1)=0$.\\
$x=y$: $x^2=9x$ $\Rightarrow$ $(0;0)$, $(9;9)$. \ $y=-1-x$: $x^2+x+5=0$ — yeçim yöq. Jami $2$ ta.
\Javob{B}

\Savol{18}
Koeffitsiyent qança katta bölsa, parabola şunça «tor»: $y=bx^2$ eng tor, $y=cx^2$ eng keng. \ $b>a>c$.
\Javob{B}

\Savol{19}
$t=x^2-5x+11$: \ $(t-1)(t+1)<48$ $\Rightarrow$ $t^2<49$ $\Rightarrow$ $t<7$ ($t>0$):
\[x^2-5x+4<0\ \Rightarrow\ 1<x<4.\quad 2+3=5.\]
\Javob{D}

\Savol{20}
\[\frac{x_1^2+x_2^2}{(x_1x_2)^2}=\frac{\frac{81}4-5}{\frac{25}4}=\frac{61}{25}.\]
\Javob{C}

\Savol{21}
$E(f)=\left[-\sqrt{52};\,\sqrt{52}\right]$, \ $7<\sqrt{52}<8$: natural sonlar $1,\dots,7$ — $7$ ta.
\Javob{B}

\Savol{22}
$5-4x-x^2=-(x-1)(x+5)$: \ $\dfrac3{5-4x-x^2}=-\dfrac12\left(\dfrac1{x-1}-\dfrac1{x+5}\right)$.
\[-\frac12\ln\frac{x-1}{x+5}\Big|_2^4=-\frac12\left(\ln\frac13-\ln\frac17\right)=-\frac12\ln\frac73=\ln\sqrt{\frac37}.\]
\Javob{B}

\Savol{23}
$u=4^x$: $f=u(2-u)^{-1/2}$, \ $f'=u'\left[(2-u)^{-1/2}+\frac u2(2-u)^{-3/2}\right]$. \ $x=0$: $u=1$, $u'=\ln4$:
\[f'(0)=\ln4\left(1+\frac12\right)=\frac32\ln4.\]
\Javob{B}

\Savol{24}
$O_1O_2=r_1+r_2=5+9=14$.
\Javob{C}

\Savol{25}
$S_{ABE}+S_{ACE}=S_{ABC}$: \ $\frac12a\cdot3+\frac12a\cdot9=\frac12a\cdot h$ $\Rightarrow$ yon tomonga balandlik $h=12$.\\
$h=a\sin30^\circ$ $\Rightarrow$ $a=24$. \ $S=\frac12\cdot24^2\cdot\sin30^\circ=144$.
\Javob{A}

\Savol{26}
$AB=15\cdot0{,}6=9$, \ $AC=12$, \ $S=\frac12\cdot9\cdot12=54$.
\Javob{B}

\Savol{27}
Kvadrat tomoni $a$; markazdan uçigaça $\frac a{\sqrt2}$, tomonigaça $\frac a2$. Har bir böyalgan uçburçak teng yonli töğri burçakli, balandligi $\frac{a\left(\sqrt2-1\right)}2$, yuzi $\frac{a^2\left(\sqrt2-1\right)^2}4$:
\[\frac{4\cdot\frac{a^2\left(3-2\sqrt2\right)}4}{a^2}=3-2\sqrt2.\]
\Javob{C}

\Savol{28}
$\tg\dfrac{180^\circ}n=\dfrac{\sqrt3}3$ $\Rightarrow$ $n=6$. \ $S=\frac12P\cdot r=\frac12\cdot12\sqrt3\cdot3=18\sqrt3$.
\Javob{A}

\Savol{29}
$\overrightarrow{AD}+\overrightarrow{DC}=\overrightarrow{AC}$, \ $\overrightarrow{CC_1}+\overrightarrow{C_1D_1}=\overrightarrow{CD_1}$. $\triangle ACD_1$ — teng tomonli (yoq diagonallari, $3\sqrt2$), vektorlar orasidagi burçak $180^\circ-60^\circ=120^\circ$:
\[3\sqrt2\cdot3\sqrt2\cdot\cos120^\circ=-9.\]
\Javob{C}

\Savol{30}
$SB^2=SA^2+AB^2$, $SD^2=SA^2+AD^2$, $SC^2=SA^2+AB^2+AD^2$ $\Rightarrow$ $SB^2+SD^2=SA^2+SC^2$:
\[36+49=SA^2+81\ \Rightarrow\ SA=2.\]
\Javob{C}

\Savol{31}
$n(A)=8$, \ $n(B)=5$, \ $n(C)=6$: \ $8+5-6=7$.
\Javob{B}

\Savol{32}
$1$ xonalidan $5$ xonaligaça: $3+9+27+81+243=363$.
\Javob{A}

\newpage
\Savol{33}
Kiçik asos $b$ (= balandlik), katta asos $2b$: \ $S_{\text{trap}}=\dfrac{b+2b}2\cdot b=\dfrac{3b^2}2$. \ Hajmlar: $\dfrac{3b^2}2H=3r^2h$ ($\pi=3$).

\textbf{33.} $H=h$: \ $\frac{3b^2}2=3r^2$ $\Rightarrow$ $\dfrac br=\sqrt2$ — \textbf{B}.

\textbf{34.} $b=2r$: \ $6r^2H=3r^2h$ $\Rightarrow$ $\dfrac hH=2$ — \textbf{F}.

\textbf{35.} $\frac{3b^2}2=6$ $\Rightarrow$ $b=2$, $r=1$: \ $\pi r^2=3$ — \textbf{C}.
\Javob{33 — B, \ 34 — F, \ 35 — C}

\newpage
\Savol{36}
$u=x+y$, $v=x-y$: \ $\dfrac4u+\dfrac4v=3$, \ $u^2+v^2=20$. \ $s=u+v$, $p=uv$:
\[4s=3p,\qquad s^2-2p=20\ \Rightarrow\ 3s^2-8s-60=0\ \Rightarrow\ s=6\ \text{yoki}\ s=-\frac{10}3.\]
$s=6$, $p=8$: $(u;v)=(2;4)$, $(4;2)$.\\
$s=-\frac{10}3$, $p=-\frac{40}9$: $t^2+\frac{10}3t-\frac{40}9=0$ — ikkita har xil işorali haqiqiy ildiz, ya'ni yana $2$ juftlik.

a) $4$ ta yeçim.

b) $x=\dfrac{u+v}2=\dfrac s2$: \ $2\cdot3+2\cdot\left(-\dfrac53\right)=\dfrac83$.
\Javob{a) 4; \quad b) $\frac83$}

\newpage
\Savol{37}
Modul içidagi ifoda manfiy emas: $\cos t=-\frac12$, bunda $t=\frac{2\pi}3\sin\frac{2x}3\in\left[-\frac{2\pi}3;\,\frac{2\pi}3\right]$. Demak $t=\pm\frac{2\pi}3$:
\[\sin\frac{2x}3=\pm1\ \Rightarrow\ \frac{2x}3=\frac\pi2+\pi k\ \Rightarrow\ x=\frac{3\pi}4+\frac{3\pi k}2.\]
$[-\pi;\,4\pi]$ da: $-\dfrac{3\pi}4,\ \dfrac{3\pi}4,\ \dfrac{9\pi}4,\ \dfrac{15\pi}4$.

a) $4$ ta. \quad b) $\dfrac{-3+3+9+15}4\pi=6\pi$.
\Javob{a) 4; \quad b) $6\pi$}

\newpage
\Savol{38}
Simmetriyadan $B(t;\,6)$, $C\left(t;\,t^2+2\right)$: \ eni $2t$, böyi $4-t^2$.
\[S(t)=2t\left(4-t^2\right)=8t-2t^3,\qquad S'=8-6t^2=0\ \Rightarrow\ t=\frac2{\sqrt3}.\]
a) $S=\dfrac4{\sqrt3}\cdot\dfrac83=\dfrac{32}{3\sqrt3}=\dfrac{32\sqrt3}9$.

b) $AC^2=\left(\dfrac4{\sqrt3}\right)^2+\left(\dfrac83\right)^2=\dfrac{48+64}9=\dfrac{112}9$ $\Rightarrow$ $AC=\dfrac{4\sqrt7}3$.
\Javob{a) $\frac{32\sqrt3}9$; \quad b) $\frac{4\sqrt7}3$}

\newpage
\Savol{39}
$x$ ni $\frac1x$ ga almaştiramiz: $\dfrac{x+1}x f\left(\frac1x\right)+f(x)=\dfrac{2x+4}{x+1}$.\\
Birinçi tenglikdan $f\left(\frac1x\right)=\dfrac{4x+2}{x+1}-(x+1)f(x)$, qöyamiz:
\[f(x)\left(1-\frac{(x+1)^2}x\right)=\frac{2x+4}{x+1}-\frac{4x+2}x\ \Rightarrow\ -\frac{x^2+x+1}x f(x)=-\frac{2\left(x^2+x+1\right)}{x(x+1)}.\]
$f(x)=\dfrac2{x+1}$.

a) $f(2)=\dfrac23$. \quad b) $f(3)=\dfrac12$, \ $f^2(3)=\dfrac14$.
\Javob{a) $\frac23$; \quad b) $\frac14$}

\newpage
\Savol{40}
Parabola öqi $x=1$, ildizi $-2$ $\Rightarrow$ ikkinçi ildiz $4$: $y_2=k(x+2)(x-4)$, $y_2(0)=-8k=8$ $\Rightarrow$ $y_2=-x^2+2x+8$.\\
Töğri çiziq $(-2;0)$, $(0;2)$ dan: $y_1=x+2$.

a) $x+2=-(x+2)(x-4)$ $\Rightarrow$ $x=3$, $y=5$: \ $x+y=8$.

b) \[S=\int_0^3\left(-x^2+2x+8-x-2\right)dx=\int_0^3\left(-x^2+x+6\right)dx=-9+\frac92+18=13{,}5.\]
\Javob{a) 8; \quad b) $13{,}5$}

\newpage
\Savol{41}
$BH^2=13^2-AH^2$, $CH^2=15^2-AH^2$, $CH-BH=\dfrac{225-169}{14}=4$ $\Rightarrow$ $BH=5$, $HC=9$, $AH=12$.

Töğri burçakli uçburçak içki aylanasi radiusi $r=\frac{a+b-c}2$:\\
$r_1=\frac{5+12-13}2=2$ $\Rightarrow$ $HM=HO=2$; \quad $r_2=\frac{9+12-15}2=3$ $\Rightarrow$ $HT=HN=3$.
\fig{\begin{tikzpicture}[scale=0.62,font=\small]
\coordinate (H) at (0,0); \coordinate (B) at (-5,0); \coordinate (C) at (9,0); \coordinate (A) at (0,12);
\coordinate (M) at (-2,0); \coordinate (O) at (0,2); \coordinate (T) at (0,3); \coordinate (N) at (3,0); \coordinate (K) at (0.5,2.5); \coordinate (P) at (0.5,0);
\fill[okgreen!30] (M)--(K)--(N)--cycle;
\draw (A)--(B)--(C)--cycle; \draw (A)--(H); \draw (-2,2) circle (2); \draw (3,3) circle (3);
\draw (M)--(K)--(N); \draw[ax] (K)--(P); \draw[ax] (K)--(0,2.5);
\node[above] at (A) {$A$}; \node[below left] at (B) {$B$}; \node[below right] at (C) {$C$}; \node[below] at (H) {$H$};
\node[below] at (M) {$M$}; \node[below] at (N) {$N$}; \node[left] at (O) {$O$}; \node[above left] at (T) {$T$}; \node[above right] at (K) {$K$}; \node[aux,below right] at (P) {$P$};
\end{tikzpicture}}
$HMO$ va $HNT$ — teng yonli töğri burçakli: $\angle KMN=\angle KNM=45^\circ$ $\Rightarrow$ $\angle MKN=90^\circ$, $MK=KN$.

a) $MN=5$: \ $S_{MKN}=\dfrac{MN^2}4=\dfrac{25}4=6{,}25$.

b) $KP$ — mediana: $P$ — $MN$ ning örtasi, $MP=2{,}5$, \ $HP=MP-MH=0{,}5$. \ $K$ dan $AH$ gaça masofa $HP=\dfrac12$.
\Javob{a) $6{,}25$; \quad b) $\frac12$}

\newpage
\Savol{42}
Perpendikulyarlar $4k$, $\sqrt5k$, $\frac8{\sqrt{21}}k$, $8\sqrt{\frac5{21}}k$ bölsin.

Aylanaga içki çizilgan törtburçakda $\triangle ANB\sim\triangle DNC$ va $\triangle AND\sim\triangle BNC$. Öxşaş uçburçaklarda mos balandliklar nisbati — öxşaşlik koeffitsiyenti:
\[k_1=\frac{AN}{DN}=\frac{BN}{CN}=\frac{4}{8/\sqrt{21}}=\frac{\sqrt{21}}2,\qquad k_2=\frac{AN}{BN}=\frac{DN}{CN}=\frac{8\sqrt5/\sqrt{21}}{\sqrt5}=\frac8{\sqrt{21}}.\]
a) $\dfrac{AN}{NC}=\dfrac{AN}{BN}\cdot\dfrac{BN}{NC}=k_2k_1=4$.

b) $AC=5NC=2\sqrt{21}$ $\Rightarrow$ $NC=\dfrac{2\sqrt{21}}5$:
\[BN=k_1NC=\frac{21}5,\qquad DN=k_2NC=\frac{16}5,\qquad BD=\frac{37}5=7{,}4.\]
\Javob{a) 4; \quad b) $7{,}4$}

\newpage
\Savol{43}
Tomon $s=\sqrt3$. $AH=GH=3s=3\sqrt3$, \ $FG=s$. Sakkizburçak yuzi — uçta ustun: $s^2+2s^2+3s^2=6s^2=18$, yarmi $9$.

$FM=m$, $AN=n$. $MGHN$ — töğri burçakli trapetsiya ($MG\parallel NH$, balandligi $GH=3\sqrt3$):
\[\frac{\left(\sqrt3-m\right)+\left(3\sqrt3-n\right)}2\cdot3\sqrt3=9\ \Rightarrow\ m+n=2\sqrt3.\]
\fig{\begin{tikzpicture}[scale=0.8,font=\small]
\coordinate (A) at (0,0); \coordinate (B) at (0,1.732); \coordinate (C) at (1.732,1.732); \coordinate (D) at (1.732,3.464); \coordinate (E) at (3.464,3.464); \coordinate (F) at (3.464,5.196); \coordinate (G) at (5.196,5.196); \coordinate (H) at (5.196,0);
\coordinate (M) at (4.878,5.196); \coordinate (N) at (2.050,0); \coordinate (Q) at (4.878,0);
\fill[okgreen!30] (M)--(G)--(H)--(N)--cycle;
\draw (A)--(B)--(C)--(D)--(E)--(F)--(G)--(H)--cycle; \draw (M)--(N); \draw[ax] (M)--(Q);
\node[below left] at (A) {$A$}; \node[left] at (B) {$B$}; \node[below right] at (C) {$C$}; \node[left] at (D) {$D$}; \node[below right] at (E) {$E$}; \node[above left] at (F) {$F$}; \node[above right] at (G) {$G$}; \node[below right] at (H) {$H$};
\node[above] at (M) {$M$}; \node[below] at (N) {$N$}; \node[aux,below] at (Q) {$Q$};
\end{tikzpicture}}
a) $P_1=5s+m+n+MN$, \ $P_2=\left(\sqrt3-m\right)+3\sqrt3+\left(3\sqrt3-n\right)+MN$:
\[P_1-P_2=5\sqrt3-7\sqrt3+2(m+n)=-2\sqrt3+4\sqrt3=2\sqrt3.\]
b) $m=\sqrt2$ $\Rightarrow$ $n=2\sqrt3-\sqrt2$. $MQ\perp AH$: $NQ=AQ-AN=\left(2\sqrt3+\sqrt2\right)-\left(2\sqrt3-\sqrt2\right)=2\sqrt2$, \ $MQ=3\sqrt3$:
\[MN=\sqrt{8+27}=\sqrt{35}.\]
\Javob{a) $2\sqrt3$; \quad b) $\sqrt{35}$}

\newpage
\Savol{44}
Silindr balandligi $h$: öq kesimida $r^2+\left(\frac h2\right)^2=R^2$ $\Rightarrow$ $h=2\sqrt{R^2-r^2}$.
\fig{\begin{tikzpicture}[scale=1.1,font=\small]
\draw (0,0) circle (1.5); \draw (-1.225,-0.866) rectangle (1.225,0.866);
\draw[axs] (0,0)--(1.225,0.866); \draw[ax] (0,0)--(1.225,0); \draw[ax] (0,-0.866)--(0,0.866);
\fill (0,0) circle (1.3pt);
\node[aux,above left] at (0.6,0.43) {$R$}; \node[aux,below] at (0.6,0) {$r$}; \node[aux,right] at (1.225,0.43) {$\frac h2$};
\end{tikzpicture}}
\[V=\pi r^2\cdot2\sqrt{R^2-r^2}.\quad x=r^2:\ \ V^2=4\pi^2x^2\left(R^2-x\right),\ \ \left(x^2R^2-x^3\right)'=2xR^2-3x^2=0\ \Rightarrow\ r^2=\frac23R^2.\]
a) $\dfrac rR=\sqrt{\dfrac23}=\dfrac{\sqrt6}3$.

b) $h=2\sqrt{\frac{R^2}3}=\frac{2R}{\sqrt3}$, \ $V_{\text{sil}}=\pi\cdot\frac23R^2\cdot\frac{2R}{\sqrt3}=\frac{4\pi R^3}{3\sqrt3}$:
\[\frac{V_{\text{şar}}}{V_{\text{sil}}}=\frac{\frac43\pi R^3}{\frac{4\pi R^3}{3\sqrt3}}=\sqrt3.\]
\Javob{a) $\frac{\sqrt6}3$; \quad b) $\sqrt3$}

\newpage
\Savol{45}
a) $A$ ga $x$, $B$ ga $24-x$ (mln söm):
\[0{,}16x+0{,}18(24-x)=4{,}09\ \Rightarrow\ 4{,}32-0{,}02x=4{,}09\ \Rightarrow\ x=11{,}5.\]
$11\,500\,000$ söm.

b) $0{,}16x=2{,}24$ $\Rightarrow$ $x=14$, \ $B$ da $10$ mln: \ $0{,}18\cdot10=1{,}8$ mln.
\Javob{a) $11\,500\,000$; \quad b) $1\,800\,000$}

\ShowJavoblar
\end{document}
